\paragraph{Unnumbered continuation: compactly supported corrections to a common fixed region.}
Retain the actual data and the exact maps of C.81, C.92 and (CFT.1):
the original closed disk $\overline D$, its interior $O$, the compatible
metric $d$, the free proper closed discrete operator group $Q$, the
boundary-fixed centralizer $\mathscr C_Q$ with a transitive actual
subgroup, the covering $\pi:O\to N=O/Q$, and the specified points
$x\in O$, $p=\pi x$. Write $\mathscr F=\mathscr C_Q(x)$.
The quotient $N$ is noncompact in this note. All compositions below
act first by the rightmost factor.

Let $(P,u_0)$ be a compact metric based space and let
$u\mapsto h_u$ be a continuous map to $\mathscr F$ in the original
uniform map-and-inverse topology, with $h_{u_0}=1$.
For every compact subpolyhedron $X\subset N$ containing $p$, the
construction below gives a based deformation in the original group,
\[
\begin{gathered}
 H:P\times[0,1]\longrightarrow\mathscr F,\qquad
 H_{u,0}=h_u,\qquad H_{u_0,t}=1,\\
 H_{u,1}|_{\pi^{-1}X}=\id,\qquad
 H_{u,t}=G_{u,t}^{-1}h_u,\qquad
 H_{u,t}^{-1}=h_u^{-1}G_{u,t},
\end{gathered}                                                    \tag{IFR.1}
\]
where $G_{u,0}=G_{u_0,t}=1$ and every $G_{u,t}^{\pm1}$ is the
identity off the inverse image of one fixed compact subset of $N$.
Every original deck label is retained. The terminal descents also
come with an explicit ordinary isotopy relative to $X$ from the
identity. This last property permits compatible successive stages
for any finite nested sequence of compact polyhedra beginning with
a closed disk neighborhood of $p$.

\emph{The exact homeomorphism-group topology downstairs.}
For a surface $M$ without boundary, the compact-open topology of
homeomorphisms equals the compact-open topology of maps and inverses.
Here is the comparison needed in this construction. Let $a\mapsto f_a$
be any family continuous for compact-open convergence of maps.
Evaluation is jointly continuous, since $M$ is locally compact.
Fix $(a_*,y_*)$ and put $z_*=f_{a_*}^{-1}(y_*)$.
Inside any prescribed neighborhood of $z_*$ choose a closed coordinate
disk $D_0$ with $z_*\in\operatorname{int}D_0$.
Choose a connected open coordinate disk $V$ about $y_*$ with compact
closure in $f_{a_*}(\operatorname{int}D_0)$.
For $a$ in a neighborhood of $a_*$, the compact-open conditions give
\[
 f_a(\partial D_0)\cap\overline V=\varnothing,
                  \qquad f_a(z_*)\in V.                         \tag{IFR.2}
\]
The two relatively open sets
$V\cap f_a(\operatorname{int}D_0)$ and $V\setminus f_a(D_0)$
partition $V$, because its points avoid $f_a(\partial D_0)$.
The first is nonempty by (IFR.2). Connectedness therefore gives
$V\subset f_a(\operatorname{int}D_0)$.
Thus $f_a^{-1}(V)\subset D_0$, proving joint inverse continuity.
The compact-open exponential law now proves continuity of inversion.
The two identity maps between the two stated topologies are continuous:
one forgets inverse coordinates, and the other uses this inversion map.
Composition is continuous by joint evaluation and the same exponential
law. These are the group topology and exact comparison used below.

Use the quotient triangulation from C.82, subdivided to make $p$
a vertex. The surface is separable: the open covering map $\pi$
sends a countable basis of $O$ to a countable basis of $N$.
Consequently the cited PL-surface results apply to this actual
surface, with its compact subpolyhedra. Its boundary is empty,
so the ``proper embedding'' condition in the source theorem imposes
no additional condition on the embeddings considered here.

\emph{The supported local extension theorem in its exact form.}
For compact subpolyhedra $K\subset X\subset M$, write
$\mathcal E_K(X,M)$ for embeddings restricting to the identity on $K$.
The primary statement
\cite[Theorem~2.1, p.~2]{C_Yagasaki2000} says that for each such
embedding $f$ and every neighborhood $U$ of $f(X)$ there are a
neighborhood $\mathcal V_f$ of $f$ and a continuous map
\[
 \phi_f:\mathcal V_f\longrightarrow
 \operatorname{Homeo}_{K\cup(M\setminus U)}(M)_0,
 \qquad \phi_f(g)f=g,qquad \phi_f(f)=\id_M.                    \tag{IFR.3}
\]
The subscript $0$ denotes the identity component. The identities
in (IFR.3) specify the actual extension and its support. The topology
comparison above supplies continuous inverses as well.

We give the finite parametric extension construction explicitly.
Let $e:P\times[0,1]\to\mathcal E_K(X,M)$ be a continuous family with
$e_{u,0}=i_X$ and $e_{u_0,t}=i_X$. This notation will be applied below
to the restrictions of a specified family of actual homeomorphisms.
The trace $\{e_{u,t}(z):u\in P,t\in[0,1],z\in X\}$ is compact.
Choose one relatively compact open $U\subset M$ containing this trace.
For every reference embedding in the compact parameter image, apply
(IFR.3) with this same $U$. The inverse images of finitely many resulting
neighborhoods cover $P\times[0,1]$.

Give this compact parameter space the maximum product metric, and
choose a Lebesgue number $\delta>0$ for that finite open cover.
Choose a time subdivision $0=a_0<\cdots<a_n=1$ with successive
lengths less than $\delta/3$. Choose finitely many points $v_i\in P$
whose open $\delta/8$ balls cover $P$, and put
\[
\begin{gathered}
 U_i=B_P(v_i,\delta/3),\qquad
 \beta_i(u)=\max\{0,\delta/4-d_P(u,v_i)\},\\
 \alpha_i(u)=\frac{\beta_i(u)}{\sum_j\beta_j(u)},\qquad
 \sum_i\alpha_i=1,qquad \operatorname{supp}\alpha_i\subset U_i.
\end{gathered}                                                    \tag{IFR.4}
\]
The denominator is positive by the chosen cover. Each support is
contained in the closed $\delta/4$ ball, which lies inside $U_i$.
For each time strip $[a,b]=[a_{k-1},a_k]$, the set
$\overline U_i\times[a,b]$ has diameter less than $\delta$.
It is therefore contained in the inverse image of one extension chart.
Denote that chart's map by $\phi_{i,k}$.

For $t\in[a,b]$ set
\[
\begin{gathered}
 t_0(u,t)=a,\qquad
 t_i(u,t)=a+(t-a)\sum_{j\leq i}\alpha_j(u),\qquad t_m(u,t)=t,\\
 \Gamma_{i,k}(u,t)=
 \phi_{i,k}(e_{u,t_i})\,
               \phi_{i,k}(e_{u,t_{i-1}})^{-1}quad(u\in U_i),\\
 \Gamma_{i,k}(u,t)=\id_M\quad(u\notin\operatorname{supp}\alpha_i),\\
 \Gamma_{i,k}(u,t)e_{u,t_{i-1}}=e_{u,t_i}.
\end{gathered}                                                    \tag{IFR.5}
\]
Here $m$ is the number of the sets $U_i$. Both arguments in the chart
formula lie in its domain. Where $\alpha_i=0$ their time arguments
coincide, making that formula the identity. The two definitions therefore
agree on the overlap of the open sets $U_i$ and
$P\setminus\operatorname{supp}\alpha_i$, which cover $P$.
This proves continuity of the pasted maps and of their explicit inverses
$\phi_{i,k}(e_{u,t_{i-1}})\phi_{i,k}(e_{u,t_i})^{-1}$.
The last identity in (IFR.5) follows by first applying the inverse
extension to $e_{u,t_{i-1}}$ to recover its reference embedding,
then applying the other extension.

Starting with $g_{u,0}=1$, define successively on the time strips
\[
\begin{gathered}
 g_{u,t}=\Gamma_{m,k}(u,t)\cdots\Gamma_{1,k}(u,t)g_{u,a},\\
 g_{u,t}^{-1}=g_{u,a}^{-1}
                   \Gamma_{1,k}(u,t)^{-1}\cdots\Gamma_{m,k}(u,t)^{-1},\\
 g_{u,t}|_X=e_{u,t},\qquad g_{u_0,t}=1,\qquad
 g_{u,t}|_{K\cup(M\setminus U)}=\id.
\end{gathered}                                                    \tag{IFR.6}
\]
At $t=a$ all factors $\Gamma_{i,k}$ are the identity, so consecutive
strip definitions agree exactly. Finite pasting proves joint continuity.
Induction using the last equation of (IFR.5) proves the restriction
identity. At $u=u_0$ both embedding arguments are the same inclusion,
so every factor is the identity, even when the chart's reference
embedding is different. The support and fixed-set assertions follow
from (IFR.3) factor by factor. Thus (IFR.6) gives the supported
extension with its inverse and based parameter behavior, without an
unwritten choice of a section of a parameter bundle.

\emph{The initial ordinary isotopy and the pointed-plane exception.}
Let $b_u=D(h_u)$, using the exact descent (CFT.1). These maps and
their inverses depend continuously on $u$, fix $p$, and belong to
the pointed identity path component by C.92.

For nontrivial $Q$, the quotient $N$ cannot be the plane. Indeed a
plane coordinate gives paths from $p$ to every point by straight
segments. Lifting those paths from $x$ gives a continuous section
$v:N\to O$ with $\pi v=\id_N$. Independence of the paths follows
by lifting their homotopies in the simply connected plane.
The image of this section is open and closed: on a connected evenly
covered neighborhood it is exactly one sheet, and all other sheets
form its open complement. Connectedness of $O$ makes it all of $O$,
so $v\pi=\id_O$ as well. Every fiber then has one point, and freeness
forces $Q=\{1\}$, contrary to the specified nontriviality.

The primary statement
\cite[Theorem~1.1, p.~1]{C_Yagasaki2000} lists the five exceptional
pairs as the plane with empty fixed set, the plane with one fixed
point, the open annulus with empty fixed set, the half-open annulus
with empty fixed set, and the punctured projective plane with empty
fixed set. For $(N,\{p\})$ and nontrivial $Q$, the empty-fixed-set
pairs do not apply and the remaining plane pair has just been excluded.
Thus the stated identity component is contractible, including every
nonorientable quotient allowed here.
Choose a contraction $C_s$ of that component, with $C_0(b)=b$ and
$C_1(b)=c$ constant. In the exact group topology established above,
\[
 B_{u,t}=C_{1-t}(b_u)\,C_{1-t}(1)^{-1},\qquad
 B_{u,0}=1,\quad B_{u,1}=b_u,\quad B_{u_0,t}=1,
                  \quad B_{u,t}(p)=p.                          \tag{IFR.7}
\]
The inverse is $C_{1-t}(1)C_{1-t}(b_u)^{-1}$. Both depend
continuously on $(u,t)$ in compact-open convergence.

For $Q=\{1\}$, the quotient map $\pi$ is the open bijection
from $O$ to $O/Q$, with its continuous inverse. We construct the
needed $B$ directly in the original disk. Choose a disk coordinate
$\iota_0:\overline D\to\overline{\mathbb D}$ and put $a=\iota_0(x)$.
The maps
\[
 m_a(z)=\frac{z-a}{1-\overline a z},\qquad
 m_a^{-1}(w)=\frac{w+a}{1+\overline a w},\qquad
 \iota=m_a\iota_0
                                                               \tag{IFR.8}
\]
are inverse disk-coordinate adjustments, with $\iota(x)=0$.
For $f_u=\iota h_u\iota^{-1}$ define
\[
 A_t(f)(z)=
 \begin{cases}
  t f(z/t),&0<t\leq1,\ |z|\leq t,\\
  z,&|z|\geq t,
 \end{cases}
 \qquad A_0(f)=\id,qquad A_t(f)^{-1}=A_t(f^{-1}).                \tag{IFR.9}
\]
The two formulas agree at $|z|=t$ because $f$ fixes the unit circle.
They give inverse homeomorphisms fixing that circle and $0$.
For $t>0$ the closed-domain pasting in $(u,t,z)$ proves joint
continuity for maps and inverses. At $t=0$ each uniform Euclidean
displacement is at most $2t$, independently of $f$, so their sum
is at most $4t$. Uniform continuity of $\iota^{-1}$ transfers
this estimate to the original metric $d$. Consequently
$\iota^{-1}A_t(f_u)\iota$ is an actual continuous based family
from $1$ to $h_u$. Descend it by the explicit maps $\pi$ and
$\pi^{-1}$ to obtain $B_{u,t}$ satisfying all of (IFR.7).
The pointed-plane exception has thus been handled by these original
boundary-fixed maps, without contracting the larger ordinary pointed
plane group.

\emph{The supported correction, its lift, and the original circle.}
Apply (IFR.3)--(IFR.6) to $M=N$, $K=\{p\}$ and
$e_{u,t}=B_{u,t}|_X$. Choose $U$ as in that construction, so
$\overline U$ is compact. Let $g_{u,t}$ be the resulting family.
Its path from the identity fixes $p$, so its unique pointed lift
from C.92 is defined. Put
\[
\begin{gathered}
 G_{u,t}|_O=\lambda(g_{u,t}),\qquad
 G_{u,t}^{-1}|_O=\lambda(g_{u,t}^{-1}),\\
 \pi G_{u,t}=g_{u,t}\pi,\qquad
 G_{u,t}(x)=x,\qquad G_{u,t}q=qG_{u,t}\quad(q\in Q),\\
 G_{u,0}=G_{u_0,t}=1,\qquad
 G_{u,t}|_{\pi^{-1}(N\setminus U)}=\id.
\end{gathered}                                                    \tag{IFR.10}
\]
The two lift formulas are inverse by uniqueness of the pointed lift
of the identity. They are jointly continuous on the interior for
maps and inverses by C.92. The same uniqueness retains each original
$q$, rather than allowing an automorphism of the deck group.
For a point over $N\setminus U$, its track lifts the constant
base track and starts at that point; it is therefore constant.
This proves the last equation, and applies to the inverse family.

Here is the quantitative original-circle extension supplied by C.81.
Choose a compact $C\subset O$ with $\pi C=\overline U$ using finitely
many local covering sections. The compact set
\[
 C'=C\cup\{G_{u,t}(z),G_{u,t}^{-1}(z):u\in P,\ t\in[0,1],\ z\in C\}
                                                               \tag{IFR.11}
\]
has one common original source bound by (C.81.8)--(C.81.9).
Explicitly, choose compact subsets $S,T$ of the given transitive
boundary-fixed centralizer subgroup, with $C\subset S(x)$ and
$C'\subset T(x)$. Set $\mathcal P=TS^{-1}\cup\{1\}$.
For $z=qc$ with $c=sx\in C$ and $G_{u,t}^{\pm1}(c)=tx\in C'$, one has
\[
 G_{u,t}^{\pm1}(z)=qtx=ts^{-1}(z),\qquad
 G_{u,t}^{\pm1}(z)\in\mathcal P(z)quad(z\in O).                \tag{IFR.12}
\]
The membership outside $\pi^{-1}\overline U$ uses the identity.
For all $a\geq0$ define
\[
 \omega_{\mathcal P}(a)=
 \max\{d(fz,fw):f\in\mathcal P,\ z,w\in\overline D,
                                                   \ d(z,w)\leq a\}.
                                                               \tag{IFR.13}
\]
Compactness makes this finite and gives $\omega_{\mathcal P}(a)\to0$
as $a\to0$. Every $f\in\mathcal P$ fixes the original circle.
Thus for every $\xi$ on that circle, (IFR.12) gives
$d(G_{u,t}^{\pm1}z,\xi)\leq\omega_{\mathcal P}(d(z,\xi))$.
Extend both maps by the identity there. This proves their joint
circle continuity uniformly in $(u,t)$, and their inverse identities
continue to the circle. Joint continuity on the compact disk is
exactly continuity of this compact parameter family in both uniform
metrics. Therefore every $G_{u,t}$ belongs to the actual $\mathscr F$.

Now define $H$ by (IFR.1). Its inverse is the displayed product;
both are continuous in the original topology, fix the original
circle and $x$, and commute with every original $q$.
To prove the endpoint assertion with its deck labels, retain the
ordinary residual family
\[
 k_{u,t}=g_{u,t}^{-1}B_{u,t},\qquad
 k_{u,t}^{-1}=B_{u,t}^{-1}g_{u,t},\qquad
 k_{u,0}=1,\qquad k_{u,t}|_X=\id.                              \tag{IFR.14}
\]
The restriction equation is exactly $g_{u,t}|_X=B_{u,t}|_X$.
For each $z\in\pi^{-1}X$, the track $\lambda(k_{u,t})(z)$
lifts the constant base track and starts at $z$; hence it stays $z$.
Uniqueness of pointed lifts gives the exact group identity
$\lambda(g^{-1}b)=\lambda(g)^{-1}\lambda(b)$.
Since $\lambda(b_u)=h_u|_O$, its value at $t=1$ is precisely
$H_{u,1}|_O$. This proves $H_{u,1}|_{\pi^{-1}X}=\id$ even when
$X$ has several components. Formula (IFR.14) also gives the claimed
ordinary isotopy relative to $X$ of the terminal descent.
The intermediate family $B$ has only been used on the specified
compact embedding trace and in this deck-label argument; its full
lift has not been substituted for the actual deformation $H$.

\emph{Exact carriers and both original-metric bounds.}
Put
\[
 K_U=\overline U,\qquad
 K_U^- =\bigcup_{u\in P}b_u^{-1}(K_U).
                                                               \tag{IFR.15}
\]
The latter is compact by joint inverse evaluation on $P\times K_U$.
The explicit products give
$H_{u,t}(z)=h_u(z)$ off $\pi^{-1}K_U^-$ and
$H_{u,t}^{-1}(z)=h_u^{-1}(z)$ off $\pi^{-1}K_U$.
These are assertions about the actual maps and their input points.

For $y=\pi z$ define the continuous compactly supported functions
\[
\begin{gathered}
 \gamma^\pm(y)=\max_{u,t}d_Q(G_{u,t}^{\pm1}z,z),\qquad
 \beta^+(y)=\max_{u,t}d_Q(G_{u,t}^{-1}h_u z,h_u z),\\
 \operatorname{supp}\gamma^\pm\subset K_U,\qquad
 \operatorname{supp}\beta^+\subset K_U^-.
\end{gathered}                                                    \tag{IFR.16}
\]
Equivariance proves independence of $z$; joint continuity and compact
maxima prove continuity, using the original topology of $d_Q$ from C.81.
The exact support statements follow from (IFR.10) and (IFR.15).
Let $a_h^\pm(y)=\max_u d_Q(h_u^{\pm1}z,z)$ be the actual envelopes
of CFT, and define for every $a\geq0$ the common inverse modulus
\[
 \Omega(a)=\max\{d(h_u^{-1}z,h_u^{-1}w):u\in P,
                \ z,w\in\overline D,\ d(z,w)\leq a\}.
                                                               \tag{IFR.17}
\]
It tends to zero at zero by compactness. Applying this same modulus
after every $q$ and using commutation gives
$d_Q(h_u^{-1}z,h_u^{-1}w)\leq\Omega(d_Q(z,w))$.
The triangle inequality now proves the simultaneous envelope estimates
\[
 a_H^+(y)\leq a_h^+(y)+\beta^+(y),\qquad
 a_H^-(y)\leq a_h^-(y)+\Omega(\gamma^+(y)),
                                                               \tag{IFR.18}
\]
where the new envelopes take the maximum over $(u,t)$.
Their continuity follows from that of the actual compact family.
The displayed bounds and compact supports prove the common $C_0$ tails.
No continuity of the monotone modulus $\Omega$ away from zero is needed.

There is also a uniform bound in the exact original sum metric.
Put $s_{u,t}=\sup_{z\in\overline D}d(G_{u,t}z,z)$.
Substitution $z=G_{u,t}w$ and symmetry of $d$ show that the same
number is the displacement supremum of $G_{u,t}^{-1}$.
Substitution $w=h_u z$ makes the forward change in (IFR.1) exactly
$s_{u,t}$; (IFR.17) bounds the inverse change. Hence
\[
 D_d(H_{u,t},h_u)\leq s_{u,t}+\Omega(s_{u,t}).                    \tag{IFR.19}
\]
All these quantities use the original closed-disk metric and the
actual lifts; no interior conjugacy is used to measure a boundary tail.

\emph{Compatible finite nested-core stages and the remaining estimate.}
Choose a finite nested chain $X_1\subset\cdots\subset X_r$ of
compact subpolyhedra, with $X_1$ a closed disk neighborhood of $p$.
The first stage above produces an actual endpoint fixing
$\pi^{-1}X_1$, and (IFR.14) places its descents in
$\operatorname{Homeo}_{X_1}(N)_0$.
For each later stage this component is contractible by the same
Theorem~1.1: its nonempty fixed set contains a disk, so none of the
five listed exceptional pairs applies, for any orientation of $N$.
Normalize that contraction by (IFR.7), now inside the group fixing
$X_i$, and use (IFR.3)--(IFR.6) with $K=X_i$ and $X=X_{i+1}$.

Each correction fixes $X_i$ throughout its path from the identity.
Its actual lift therefore fixes every point of $\pi^{-1}X_i$
throughout, by the same constant-track argument. All earlier fixed
points stay fixed during the next actual deformation. Equation
(IFR.14) supplies the relative component witness for the next endpoint.
Induction gives all stages and their explicitly inverse maps; finite
concatenation is continuous in the original topology at the matching
endpoints. This construction includes compact polyhedra with edges,
cycles and complementary components through their actual embeddings.

For each finite stage, (IFR.15)--(IFR.19) prove its carriers and its
original map-and-inverse control. They do not bound $s_{u,t}$ by the
initial displacement envelope outside the already fixed core.
In particular no estimate forcing these correction sizes to tend to
zero along an exhaustion has been proved here. The remaining global
task is that quantitative estimate for a compatible infinite sequence
in the original topology. The present result is the constructed
common fixed region and the compatible finite stages just proved.
